\section{Kuantor Universal: Properti Global dan Implikasi}

Kuantor universal digunakan untuk menyatakan bahwa suatu sifat berlaku untuk semua elemen dalam domain \cite{rosen2019}.

\subsection{Definisi Formal Kuantor Universal}

\begin{definisi}[Kuantor Universal]
Kuantifikasi universal merubah fungsi proposisional $P(x)$ menjadi proposisi ``$P(x)$ bernilai benar untuk semua nilai $x$ di dalam semesta pembicaraan $\mathcal{U}$''. Notasinya ditulis:
\[ \forall x \, P(x) \]
Simbol $\forall$ asalnya diambil dari huruf ``A'' tebalik (\textit{All}).
\end{definisi}

Dibaca:
\begin{itemize}
    \item ``Untuk setiap $x$, $P(x)$''
    \item ``Untuk semua $x$, berlaku $P(x)$''
    \item ``Berapapun $x$ yang dipilih, $P(x)$ benar''
\end{itemize}

Jika semesta pembicaraan $\mathcal{U}$ berbentuk himpunan terhingga $\mathcal{U} = \{a_1, a_2, \ldots, a_n\}$, maka kuantor universal ekuivalen secara logis dengan konjungsi seluruh elemen domain:
\[ \forall x \, P(x) \equiv P(a_1) \land P(a_2) \land \ldots \land P(a_n) \]

Karena ekuivalen dengan konjungsi, $\forall x \, P(x)$ bernilai \textbf{Benar} jika dan hanya jika $P(x)$ benar untuk \emph{setiap} elemen. Jika ditemukan satu elemen $c$ dengan $P(c)$ bernilai salah (disebut \textit{counterexample} atau contoh penyangkal), klaim universal bernilai salah.

\subsection{Kombinasi dengan Implikasi ($\rightarrow$)}

Dalam translasi dari bahasa alami, kuantor universal sangat erat kaitannya dengan operator implikasi. Pernyataan kategoris seperti "Semua A adalah B" selalu ditranslasikan menjadi struktur implikasi.

\begin{contoh}
Translasi proposisi: ``Setiap mahasiswa mematuhi peraturan.''
\begin{itemize}
    \item \textbf{Opsi 1: Domain terbatas}. Jika semesta $\mathcal{U}$ adalah himpunan seluruh mahasiswa, dengan $P(x) \equiv \text{``}x \text{ mematuhi peraturan''}$, notasinya:
    \[ \forall x \, P(x) \]
    
    \item \textbf{Opsi 2: Domain luas}. Jika semesta $\mathcal{U}$ adalah himpunan seluruh makhluk hidup, subjek mahasiswa perlu dinyatakan eksplisit. Misal $M(x) \equiv \text{``}x \text{ mahasiswa''}$. Notasinya:
    \[ \forall x \, (M(x) \rightarrow P(x)) \]
\end{itemize}
\end{contoh}

\begin{catatan}
\textbf{Perhatian!} Kuantor universal umumnya \textbf{tidak} dipasangkan dengan konjungsi ($\land$) untuk pola ``setiap A adalah B''. Bentuk $\forall x \, (M(x) \land P(x))$ berarti seluruh entitas di semesta adalah mahasiswa dan sekaligus mematuhi peraturan. Untuk pola kategoris ``setiap A adalah B'', bentuk yang tepat adalah implikasi: $\forall x \, (M(x) \rightarrow P(x))$.
\end{catatan}

\subsection{Contoh Matematis Universal}

\begin{contoh}
Diberikan domain $\mathcal{U} = \text{Bilangan Real } (\mathbb{R})$.
\begin{enumerate}
    \item $\forall x \ (x^2 \geq 0)$ \\ \textbf{Nilai: Benar}. Kuadrat setiap bilangan riil tidak kurang dari 0.
    \item $\forall x \ (x + 1 > x)$ \\ \textbf{Nilai: Benar}. Menambahkan 1 selalu menghasilkan nilai yang lebih besar.
    \item $\forall x \ (x^2 > x)$ \\ \textbf{Nilai: Salah}. Untuk $x = 0.5$, diperoleh $0.5^2 = 0.25$ dan $0.25 \not> 0.5$. Nilai $x=0.5$ adalah \textit{counterexample}.
\end{enumerate}
\end{contoh}
